\documentclass[12pt]{article}
\usepackage[margin=1in]{geometry}
\usepackage[T1]{fontenc}
\usepackage{lmodern}
\usepackage{array}
\usepackage{longtable}
\usepackage{hyperref}
\usepackage{graphicx}
\usepackage{amsmath}
\usepackage{amssymb}

\title{Project Report: Multi-Packet GMSK Communication System with CSV Interface}
\author{}
\date{\today}

\begin{document}
\maketitle

\section{Introduction}
This report documents a complete MATLAB implementation of a GMSK (Gaussian Minimum Shift Keying) communication system supporting multi-packet transmission with AXI4-Stream CSV interface. The system implements a full digital communication pipeline including packet building, convolutional encoding, GMSK modulation, AWGN channel simulation, preamble detection, soft-decision Viterbi decoding, and CSV-based data input/output.

The system is designed to transmit arbitrary-length data payloads across multiple packets, with each packet containing a 34-byte information payload encoded using a rate-1/2 convolutional code. The receiver performs correlation-based preamble detection to identify packet boundaries and recovers the transmitted data through soft-decision Viterbi decoding.

\section{System Overview}

\subsection{Architecture}
The system follows a conventional digital communication architecture:

\begin{figure}[h]
\centering
\begin{verbatim}
CSV Input → Packet Builder → Conv Encoder → GMSK Modulator → Channel
                                                                 ↓
CSV Output ← Byte Conversion ← Viterbi Decoder ← Payload Extract ← Preamble Detect ← GMSK Demodulator
\end{verbatim}
\caption{System Block Diagram}
\end{figure}

\subsection{Key Features}
\begin{itemize}
    \item Support for arbitrary-length payload transmission via multi-packet aggregation
    \item Rate-1/2 convolutional encoding with K=7 constraint length
    \item Gaussian pulse shaping for spectral efficiency
    \item Soft-decision Viterbi decoding for improved error correction
    \item Correlation-based preamble detection for packet synchronization
    \item AXI4-Stream CSV format for data input/output
    \item Comprehensive diagnostic plots and error metrics
\end{itemize}

\section{Packet Structure}

The Waveform class (waveformId=0) defines a fixed packet structure optimized for robust communication:

\subsection{Packet Composition}
Each packet contains the following fields:

\begin{longtable}{|c|c|p{8cm}|}
\hline
\textbf{Field} & \textbf{Length (bits)} & \textbf{Description} \\
\hline
\endfirsthead
\hline
\textbf{Field} & \textbf{Length (bits)} & \textbf{Description} \\
\hline
\endhead
\hline
\endfoot
\hline
Preamble & 48 & Known synchronization sequence for detection \\
\hline
Payload Chunk 1 & 136 & First quarter of encoded payload (136 bits) \\
\hline
Training Sequence 1 & 8 & First quarter of training sequence \\
\hline
Payload Chunk 2 & 136 & Second quarter of encoded payload \\
\hline
Training Sequence 2 & 8 & Second quarter of training sequence \\
\hline
Payload Chunk 3 & 136 & Third quarter of encoded payload \\
\hline
Training Sequence 3 & 8 & Third quarter of training sequence \\
\hline
Payload Chunk 4 & 136 & Fourth quarter of encoded payload \\
\hline
Training Sequence 4 & 8 & Fourth quarter of training sequence \\
\hline
\textbf{Total Burst} & \textbf{624} & Complete packet length \\
\hline
\end{longtable}

\subsection{Payload Parameters}
\begin{itemize}
    \item Information payload: 34 bytes (272 bits)
    \item Encoded payload: 544 bits (rate 1/2 convolutional coding)
    \item Preamble: 48 bits (pseudo-random sequence for correlation detection)
    \item Training sequences: 32 bits total (8 bits $\times$ 4, interleaved)
    \item Total burst: 624 bits
\end{itemize}

The interleaved structure places training sequences between payload chunks to assist with channel estimation and provide known reference points within the burst.

\section{Transmitter Components}

\subsection{Waveform Class (\texttt{Waveform.m})}
The Waveform class defines the packet structure parameters:
\begin{itemize}
    \item \texttt{waveformId}: Identifies the waveform configuration (0 for FEC rate 1/2)
    \item \texttt{pldLenInfBytes}: Information payload length in bytes (34)
    \item \texttt{pldLenEncBits}: Encoded payload length in bits (544)
    \item \texttt{prmbSeqLen}: Preamble sequence length (48 bits)
    \item \texttt{trngSeqLen}: Training sequence length (32 bits)
    \item \texttt{burstBlockLen}: Total burst length (624 bits)
\end{itemize}

The class generates reproducible preamble and training sequences using a fixed random seed (12345) for consistency between transmission and reception.

\subsection{Convolutional Encoder (\texttt{convEncoder.m})}
The convolutional encoder implements a rate-1/2 code with K=7 constraint length:
\begin{itemize}
    \item Generator polynomials: [171 133] (octal)
    \item Constraint length: K = 7
    \item Code rate: 1/2
    \item Input: 272 information bits
    \item Output: 544 encoded bits
\end{itemize}

The encoder uses a shift register implementation where each input bit produces two output bits by computing the parity of specific tap combinations.

\subsection{Gaussian Filter (\texttt{GaussianFilter.m})}
The Gaussian filter performs pulse shaping to limit spectral occupancy:
\begin{itemize}
    \item BT product: 0.35 (typical for GMSK)
    \item Samples per symbol (SPS): 4
    \item Filter span: 8 symbols
\end{itemize}

The Gaussian filter shape is characterized by the BT (bandwidth-time) product, which determines the tradeoff between spectral efficiency and intersymbol interference.

\subsection{GMSK Mapper (\texttt{GmskMapper.m})}
The GMSK mapper converts filtered bits to complex baseband signals:
\begin{itemize}
    \item LUT size: 512 entries
    \item Accumulator bits: 16
    \item Bit gain reduction: 8
    \item Quantization: Q15 format
\end{itemize}

The mapper implements phase accumulation where each bit advances the carrier phase by $\pi/2$ radians, producing the characteristic GMSK constant-envelope signal.

\section{Receiver Components}

\subsection{RxDemodulator (\texttt{RxDemodulator.m})}
The receiver demodulator converts the received complex signal to soft metrics:
\begin{itemize}
    \item Samples per symbol: 4 (must match transmitter)
    \item Gaussian filter span: 8 symbols
    \item Output: Log-likelihood ratios (LLRs) for each bit
\end{itemize}

The demodulator applies matched filtering and samples at symbol intervals, accounting for the group delay introduced by the transmitter filter.

\subsection{GMSK Demapper (\texttt{GmskDemapper.m})}
The demapper extracts phase information from the received signal:
\begin{itemize}
    \item Computes phase difference between consecutive samples
    \item Unwraps phase to recover continuous phase trajectory
    \item Outputs soft metrics proportional to bit reliability
\end{itemize}

Phase-based demodulation is optimal for GMSK since the information is encoded in phase changes.

\subsection{Viterbi Decoder (\texttt{ViterbiDecoder.m})}
The Viterbi decoder performs soft-decision convolutional decoding:
\begin{itemize}
    \item Constraint length: K = 7
    \item Generator polynomials: [171 133] octal
    \item Traceback length: 35
    \item Input type: Soft LLR
    \item Soft bit width: 3 bits
    \item Mode: Truncated
\end{itemize}

The decoder converts LLRs to soft metrics using the convention that metric 0 indicates strong '0' and metric $2^N-1$ indicates strong '1'.

Key parameters:
\begin{itemize}
    \item \texttt{nVar}: Noise variance for LLR scaling
    \item \texttt{llrClip}: LLR clipping range (default 127)
    \item \texttt{tracebackLen}: Number of trellis steps to trace back
\end{itemize}

\section{Packet Detection and Processing}

\subsection{Preamble Detection (\texttt{detect\_preamble.m})}
The preamble detector uses normalized cross-correlation:
\begin{enumerate}
    \item Compute correlation between received LLR stream and known preamble sequence
    \item Normalize correlation by preamble length
    \item Apply threshold (typically 0.6) to identify peaks
    \item Apply non-maximum suppression with minimum spacing (e.g., 500 bits)
    \item Return ordered list of detections with indices and correlation values
\end{enumerate}

The detection threshold controls the tradeoff between sensitivity and false alarm rate. Lower thresholds detect more packets but may produce false positives.

\subsection{Packet Building (\texttt{build\_packet.m})}
The packet builder constructs complete bursts:
\begin{enumerate}
    \item Convert payload bytes to bits (MSB-first)
    \item Apply convolutional encoding (272 $\rightarrow$ 544 bits)
    \item Insert preamble at the start
    \item Interleave payload chunks with training sequences
    \item Output complete 624-bit burst
\end{enumerate}

The interleaved structure ensures training sequences are distributed throughout the packet for robust channel estimation.

\subsection{Payload Extraction (\texttt{extract\_payload.m})}
The payload extractor removes overhead:
\begin{enumerate}
    \item Skip preamble bits
    \item Extract 4 payload chunks, discarding training sequences
    \item Concatenate chunks into continuous encoded payload
    \item Output 544 encoded bits
\end{enumerate}

\section{Multi-Packet Processing}

The system supports transmission of arbitrary-length data by splitting across multiple packets.

\subsection{Multi-Packet Builder (\texttt{build\_multipacket.m})}
\begin{enumerate}
    \item Read all input bytes from CSV
    \item Split into 34-byte chunks (last chunk may be padded)
    \item Build each packet using \texttt{build\_packet}
    \item Insert configurable guard bits between packets
    \item Return concatenated bitstream with packet metadata
\end{enumerate}

Default guard bits: 100 bits between packets to ensure separation.

\subsection{Multi-Packet Reception (\texttt{receive\_multipacket.m})}
\begin{enumerate}
    \item For each detection from preamble detector:
    \begin{itemize}
        \item Extract burst from LLR stream
        \item Remove preamble and training sequences
        \item Apply Viterbi decoding
        \item Convert bits to bytes
    \end{itemize}
    \item Concatenate all recovered packet bytes
    \item Trim padding from final packet
    \item Return aggregate bytes with per-packet status
\end{enumerate}

The receiver handles partial packet failures gracefully, continuing to process detected packets even if some fail.

\section{CSV Interface}

The system uses AXI4-Stream format for CSV input/output:

\subsection{Input Format (\texttt{read\_axi\_stream\_csv.m})}
\begin{verbatim}
TDATA,TKEEP,TLAST
9B553F84,f,0
0789DC7F,f,0
CA04CA04,3,1
\end{verbatim}

\begin{itemize}
    \item TDATA: 32-bit word in hexadecimal (8 hex digits)
    \item TKEEP: Byte enable mask (f = all 4 bytes valid, 3 = upper 2 bytes)
    \item TLAST: End of frame flag (1 on last word)
\end{itemize}

Byte ordering: TDATA stores bytes in big-endian format [B3 B2 B1 B0] where B3 is MSB.

\subsection{Output Format (\texttt{write\_axi\_stream\_csv.m})}
Same format as input, with decoded bytes written in TDATA field.

\section{Test Scripts and Validation}

\subsection{Single-Packet Test (\texttt{test\_complete\_system.m})}
Tests end-to-end transmission of a single 34-byte packet:
\begin{itemize}
    \item Eb/N0: 10 dB
    \item SPS: 4
    \item BT: 0.35
    \item Preamble threshold: 0.65
    \item Generates 9 diagnostic plots including constellation, correlation, and error analysis
\end{itemize}

Output metrics:
\begin{itemize}
    \item Byte Error Rate (BER)
    \item Bit Error Rate
    \item Error locations
    \item Transmitted vs received byte comparison
\end{itemize}

\subsection{Multi-Packet Test (\texttt{test\_multipacket\_system.m})}
Tests continuous multi-packet transmission:
\begin{itemize}
    \item Eb/N0: 30 dB
    \item SPS: 4
    \item BT: 0.35
    \item Preamble threshold: 0.6
    \item Guard bits: 100
    \item Minimum packet spacing: 500 bits
    \item Generates 12 diagnostic plots
\end{itemize}

Output metrics:
\begin{itemize}
    \item Number of transmitted/detected/recovered packets
    \item Packet success rate
    \item Byte Error Rate
    \item Bit Error Rate
    \item Per-packet status
    \item Detection timeline
\end{itemize}

\subsection{Simple Example (\texttt{simple\_multipacket\_example.m})}
Lightweight demonstration script for quick testing.

\section{Performance Parameters}

\subsection{System Parameters}
\begin{longtable}{|c|c|p{8cm}|}
\hline
\textbf{Parameter} & \textbf{Typical Value} & \textbf{Description} \\
\hline
\endfirsthead
\hline
\textbf{Parameter} & \textbf{Typical Value} & \textbf{Description} \\
\hline
\endhead
\hline
\endfoot
\hline
Eb/N0 & 10-30 dB & Energy per bit to noise spectral density \\
\hline
SPS & 4 & Samples per symbol \\
\hline
BT & 0.35 & Gaussian filter bandwidth-time product \\
\hline
Preamable Threshold & 0.6-0.65 & Correlation detection threshold \\
\hline
Guard Bits & 100 & Bits between packets in multi-packet mode \\
\hline
Min Packet Spacing & 500 bits & Minimum detection spacing for NMS \\
\hline
\end{longtable}

\subsection{Expected Performance}
\begin{longtable}{|c|c|c|}
\hline
\textbf{Eb/N0 (dB)} & \textbf{Expected BER} & \textbf{Comments} \\
\hline
\endfirsthead
\hline
\textbf{Eb/N0 (dB)} & \textbf{Expected BER} & \textbf{Comments} \\
\hline
\endhead
\hline
\endfoot
\hline
3 & ~0.05-0.1 & Marginal performance \\
\hline
10 & ~0.001-0.01 & Usable for most applications \\
\hline
30 & ~0.0001 & Excellent performance, near error-free \\
\hline
\end{longtable}

\section{File Inventory}

\subsection{Core Communication Blocks}
\begin{longtable}{|p{4cm}|p{10cm}|}
\hline
\textbf{File} & \textbf{Description} \\
\hline
\endfirsthead
\hline
\textbf{File} & \textbf{Description} \\
\hline
\endhead
\hline
\endfoot
\hline
\texttt{convEncoder.m} & Rate 1/2 convolutional encoder (K=7) \\
\hline
\texttt{ViterbiDecoder.m} & Soft-decision Viterbi decoder \\
\hline
\texttt{GaussianFilter.m} & Gaussian pulse shaping filter \\
\hline
\texttt{GmskMapper.m} & GMSK modulator with LUT \\
\hline
\texttt{GmskDemapper.m} & Phase detection and unwrapping \\
\hline
\texttt{RxDemodulator.m} & Complete demodulator with LLR output \\
\hline
\end{longtable}

\subsection{Packet Processing}
\begin{longtable}{|p{4cm}|p{10cm}|}
\hline
\textbf{File} & \textbf{Description} \\
\hline
\endfirsthead
\hline
\textbf{File} & \textbf{Description} \\
\hline
\endhead
\hline
\endfoot
\hline
\texttt{Waveform.m} & Defines packet structure parameters \\
\hline
\texttt{build\_packet.m} & Constructs single packet from payload \\
\hline
\texttt{build\_multipacket.m} & Builds multi-packet transmission stream \\
\hline
\texttt{extract\_payload.m} & Extracts payload from received burst \\
\hline
\texttt{detect\_preamble.m} & Correlation-based preamble detection \\
\hline
\texttt{receive\_multipacket.m} & Processes multiple detected packets \\
\hline
\end{longtable}

\subsection{CSV Interface}
\begin{longtable}{|p{4cm}|p{10cm}|}
\hline
\textbf{File} & \textbf{Description} \\
\hline
\endfirsthead
\hline
\textbf{File} & \textbf{Description} \\
\hline
\endhead
\hline
\endfoot
\hline
\texttt{read\_axi\_stream\_csv.m} & Reads AXI4-Stream format CSV \\
\hline
\texttt{write\_axi\_stream\_csv.m} & Writes AXI4-Stream format CSV \\
\hline
\end{longtable}

\subsection{Test Scripts}
\begin{longtable}{|p{4cm}|p{10cm}|}
\hline
\textbf{File} & \textbf{Description} \\
\hline
\endfirsthead
\hline
\textbf{File} & \textbf{Description} \\
\hline
\endhead
\hline
\endfoot
\hline
\texttt{test\_complete\_system.m} & Single-packet end-to-end test \\
\hline
\texttt{test\_multipacket\_system.m} & Multi-packet end-to-end test \\
\hline
\texttt{simple\_multipacket\_example.m} & Lightweight demonstration \\
\hline
\texttt{compare\_multipacket\_results.m} & Diagnostic visualization utility \\
\hline
\end{longtable}

\section{Usage Example}

\subsection{Running the System}
\begin{verbatim}
% Single packet test
test_complete_system

% Multi-packet test
test_multipacket_system
\end{verbatim}

These scripts will:
\begin{enumerate}
    \item Read input data from CSV (tx\_data.csv or tx\_multipacket\_data.csv)
    \item Build and modulate packets
    \item Add AWGN channel noise
    \item Demodulate and detect preambles
    \item Decode payloads with Viterbi algorithm
    \item Write recovered data to output CSV
    \item Display diagnostic plots and error statistics
\end{enumerate}

\section{Conclusion}

This project implements a complete GMSK communication system with robust error correction through convolutional coding and Viterbi decoding. The multi-packet support enables transmission of arbitrary-length data streams, while the CSV interface provides simple integration with external systems using the AXI4-Stream standard.

Key achievements:
\begin{itemize}
    \item Complete transmitter and receiver chain implementation
    \item Support for arbitrary-length payload transmission
    \item Comprehensive diagnostic and visualization tools
    \item AXI4-Stream CSV format for easy data input/output
    \item Configurable parameters for different operating conditions
\end{itemize}

The system demonstrates reliable communication at moderate to high SNR levels and provides a solid foundation for further development such as higher-order modulation, channel estimation improvements, or hardware implementation.

\end{document}
